\documentclass[11pt,a4paper]{article}
\usepackage[T1]{fontenc}
\usepackage{lmodern,amsmath,amssymb,booktabs,microtype}
\usepackage[margin=24mm]{geometry}
\usepackage[colorlinks=true,urlcolor=blue,linkcolor=blue]{hyperref}
\newcommand{\EE}{\mathbb E}
\newcommand{\tr}{\operatorname{tr}}
\newcommand{\W}[1]{\mathsf W_{#1}}
\newcommand{\Db}{\overline D}
\newcommand{\J}{J_s}
\title{Two orientations of the fermionic force\\The projected second memory coefficient in energy space}
\author{Alan Day \and Chengheng\thanks{AI research assistant (OpenAI).}}
\date{Local research note --- 6 November 2026}
\begin{document}\maketitle
\begin{abstract}
The output-rank deficit in the second memory coefficient is asymptotic
to the average of the annihilation and creation force intensities.
We prove this reduction using an exact Casimir identity and overlap
bounds. The two force orientations have a positive matrix of
$q^2$ probe measures. An endpoint certificate proves that their
cross measures with the fourth-chaos residual vanish at every fixed
energy order. The second memory coefficient, after the previously
specified eight-direction projection, therefore has a positive sum
of the existing $q^4$ and $q^2$ Mehler kernels.
\end{abstract}
\section{State, conventions, and result}
Use standard CAR, the full overlapping complex Gaussian $p$-body
Hamiltonian, fixed charge $m=N/2$, and the normalized maximally mixed
trace. All inner products are annealed:
$\langle X,Y\rangle=\EE\tr_m X^\dagger Y$. Set $\hbar=1$ and
\begin{equation}
 h=H_m,\quad \lambda_N=p^2/N\longrightarrow\lambda>0,\quad
 q=e^{-4\lambda},\quad C_N=\EE\tr_mh^2\longrightarrow C=\J^2/\lambda.
\end{equation}
The variance convention and the same-coupling charge lifts are those
of the preceding notes. $\W j$ denotes ordinary Gaussian-coordinate
chaos, not the spectral $q$-Wick polynomial of Ref.~\cite{TT}.
Define $d_i=[H,c_i]$ on full Fock space and restrict neutral operators
to charge $m$:
\begin{equation}
 A_- =\sum_i d_i^\dagger d_i,\qquad A_+=\sum_i d_id_i^\dagger,
 \quad O_\pm=\W2 A_\pm,\quad B=\W2h^2.
\end{equation}
Let $\widehat\Pi_N$ be the actual finite-$N$ orthogonal projection onto
$\operatorname{span}\{I,B,O_-,J_hF,E_0,M_0,S_0,Z_0\}$, as in
Ref.~\cite{desc}. For $X_N=(1-\widehat\Pi_N)mZ_2$, our result is
\begin{equation}
 \boxed{\quad d\nu_X=V_4\,d\mu_{q^4}+V_2\,d\mu_{q^2},\qquad
 C_X(t)=V_4G_{q^4}(t)+V_2G_{q^2}(t).\quad} \label{result}
\end{equation}
Here $\mu_r$ is the normalized, existing two-energy probe measure,
$G_r(0)=1$, and both weights are positive. This is the full
two-point spectrum of this specified projected coefficient; the
finite-time memory operator remains a separate problem.

\newpage
\section{An exact bridge from output rank to force orientation}
Let $Q=\sum_i c_i^\dagger c_i$ and define on neutral full-Fock operators
\begin{equation}
 DX=QX-\sum_i c_i^\dagger Xc_i,\qquad
 \Db X=(N-Q)X-\sum_i c_iXc_i^\dagger.
\end{equation}
For the neutral quadratic Casimir $\mathcal C_x$, with eigenvalues
$x_a=a(N+1-a)$ on rank $a$, CAR gives at half filling
\begin{equation}
 \mathcal C_xX=(m+1)DX+m\Db X-D\Db X,\qquad DH=pH.
 \label{casimir}
\end{equation}
The compositions on the right are formed before charge restriction.
Write $h=\sum_a h_a$, $u=p-a$, and
\begin{equation}
 D_{\rm in}h=\sum_a u h_a,\qquad
 E_{\rm in}=\sum_a\frac{u(u-1)}{m-p+1}h_a.
\end{equation}
Equation~\eqref{casimir} implies the exact relation
\begin{equation}
 (p-\Db)h=2D_{\rm in}h+E_{\rm in}.
 \label{input}
\end{equation}
Indeed $\Db h_a=(x_a-pm)h_a/(m-p+1)$, and
$(x_p-x_a)/(m-p+1)=2u+u(u-1)/(m-p+1)$.

Let $Z^+_{ab;r}$ be the rank-$r$ component of
$h_ah_b+h_bh_a$ for $a<b$, and of $h_a^2$ for $a=b$.
Define the actual deficit and its Casimir proxy by
\begin{align}
 \mathcal S_N&=\sum_{a\leq b,r}(a+b-r)\W2Z^+_{ab;r},\\
 \widetilde{\mathcal S}_N
 &=\sum_{a\leq b,r}\frac{x_{a+b}-x_r}{N}\W2Z^+_{ab;r}.
\end{align}
The earlier positive rank-moment bounds give
$\|\mathcal S_N-\widetilde{\mathcal S}_N\|_2=O(p^{-1})$.
Introduce full-Fock operators and then restrict:
\begin{align}
 \overline A&=(2p-\Db)H^2
       =A_++\{H,(p-\Db)H\},\\
 E_{\rm dbl}&=(2p-D)\overline A,\qquad
 U_{{\rm in},2}=\{h,D_{\rm in}^2h\}+2(D_{\rm in}h)^2.
\end{align}
Substitution into~\eqref{casimir}, with
$A_-=(2p-D)H^2$, gives an exact finite-$N$ identity:
\begin{align}
 N\left(\widetilde{\mathcal S}_N-\frac{O_-+O_+}{2}\right)
 ={}&(1-2p)O_--2pO_+-\W2\{h,D_{\rm in}h\}\notag\\
 &-\W2 U_{{\rm in},2}+\W2 E_{\rm dbl}
 +(m-2p)\W2\{h,E_{\rm in}\}. \label{bridge}
\end{align}
It follows either by collecting $x_a+x_b-2ab-x_r$, or by applying
the Casimir directly to $h^2$. No division by a possibly vanishing
rank product is used.

\newpage
\section{Controlling the bridge error}
The only new bound needed in~\eqref{bridge} is the normal/reverse
double mark. We give its local reason explicitly. On a neutral word,
the sitewise maps in $D=\sum_iD_i$, $\Db=\sum_i\Db_i$ have the table
below; fermion signs in odd local factors are fixed by the original
global ordering. Put $n=c^\dagger c$.
\begin{center}
\begin{tabular}{c|ccccc}
local factor&$I$&$c$&$c^\dagger$&$n$&$1-n$\\\hline
$D_i$&$0$&$0$&$c^\dagger$&$n$&$-n$\\
$\Db_i$&$0$&$c$&$0$&$n-1$&$1-n$
\end{tabular}
\end{center}
Consider one two-slot monomial in $H^2$, before summing indices.
Let $a_i,b_i$ count its original creation and annihilation factors
at site $i$. Then $\sum_i a_i=\sum_i b_i=2p$ and
\begin{equation}
 (2p-D)(2p-\Db)=
 \left[\sum_i(a_i-D_i)\right]\left[\sum_j(b_j-\Db_j)\right].
\end{equation}
At a site used by just one of the four original index sets, both
deficits annihilate its local factor. The remaining sites number
at most $K$, the sum of pairwise index-set overlaps. Each local
map has a bounded coefficient norm. Thus one deficit has coefficient
bound $K_0K$, and two deficits have bound $K_0^2K^2$ on the overlap
sites. A prior deficit does not modify a singleton site, so the
second operation creates no hidden factor of $p$ there.

In the Gaussian second chaos, internal contractions of a two-slot
block are absent. For each fixed Wick graph, its distinct $p$-sets
have uniformly bounded $\EE e^{aK}$ for every fixed $a$ at fixed
$p^2/N$. This follows from the hypergeometric factorial-moment bound
for each pair and H\"older for the finite number of pairs. After local
occupation expansion, absolute trace coefficients are bounded by
$K_1(1+K)^j e^{K_2K}$ times the singleton normalization.
The same argument treats any fixed number of blocks and marks.

For the half-filled trace, the exact local/full-Fock ratio is
\begin{equation}
 R_N(L,t)=2^L\frac{(m)_t^\downarrow(m)_{L-t}^\downarrow}
 {(N)_L^\downarrow}\leq e^{K_3p^2/N}.
\end{equation}
On fixed overlap windows the balanced singletons give
$|2t-L|\leq K_4K$, hence $R_N\to1$. The exponential overlap
bound controls the complementary window, as in Ref.~\cite{desc}.
Consequently $\|O_\pm\|_2$ and $\|\W2E_{\rm dbl}\|_2$ are bounded.
The earlier factorial input-rank tails also bound every fixed
moment of $D_{\rm in}^j h$, while those of $E_{\rm in}$ are
$O(N^{-1})$ in norm. Gaussian fourth-moment bounds and H\"older
therefore bound the products in~\eqref{bridge}. We obtain
\begin{equation}
 \boxed{\ \mathcal S_N=\tfrac12(O_-+O_+)+O_{L^2}(p^{-1}).\ } \label{strong}
\end{equation}
The limits and constants here keep $\lambda>0$ fixed.

\newpage
\section{The oriented two-energy spectral matrix}
For two marked pairs $(a,b)$ and $(c,d)$ in a Wick word, let
$\pi$ be the Gaussian pairing. The local determinant calculation
gives the same single-mark coefficient for either orientation,
$\alpha=2-4\,\mathbf1_{\rm incident\ crossing}$. The mixed coefficient is
\begin{equation}
 \delta=\begin{cases}
 4\,\mathbf1_{\pi(a)=d,\,\pi(b)=c},&\text{equal orientations},\\
 -4\,\mathbf1_{\pi(a)=c,\,\pi(b)=d},&\text{opposite orientations}.
 \end{cases} \label{minor}
\end{equation}
The determinant entry for a normal mark $(x,y)$ changes by $-2$
at $(\pi(j),i)=(x,y)$; for a reverse mark it changes by $+2$
at $(\pi(j),i)=(y,x)$. These statements hold before any spectrum
is postulated. The two-mark weight is
$q^{\rm cr}(\lambda\delta+\lambda^2\alpha\beta)$ times the
appropriate power of $C$.

For arbitrary nonnegative $a,b$, put $n=a+b$. In unit-variance
energy variables let $P_{ab}$ be the mixed moment of $\mu_{q^2}$,
let $m_j$ be the $q$-Gaussian moment, and put
\begin{equation}
 u_q=\frac{1-q}{1+q},\qquad S_n=m_{n+4}-2m_{n+2}+m_n.
\end{equation}
Swapping the two adjacent endpoints of each force insertion changes
only their mutual crossing. The usual four-member exchange orbits
give $u_q^2$; the completely connected double bridge has two
members. For opposite orientations only its crossing member
contributes to $\delta$. Thus for all two-arc words
\begin{align}
 \sum q^{\rm cr}\delta_{--}&=4P_{ab},&
 \sum q^{\rm cr}\delta_{-+}&=-4qP_{ab},\\
 \sum q^{\rm cr}\alpha\beta&=4u_q^2S_n+
                  \frac{16q}{1+q}P_{ab},&
 \sum q^{\rm cr}\alpha&=2u_qS_n.
\end{align}
The last identity also fixes both crosses with the conserved $B$.
Set $c=2\lambda u_q$ and $Y_\pm=O_\pm-cB$. Subtracting this common
background yields the two-energy matrix
\begin{equation}
 d\nu_{(Y_-,Y_+)}=
 \begin{pmatrix}w&x\\x&w\end{pmatrix}d\mu_{q^2},\quad
 \begin{aligned}
 w&=4C^2\lambda+\frac{16C^2\lambda^2q}{1+q},\\
 x&=-4C^2\lambda q+\frac{16C^2\lambda^2q}{1+q}.
 \end{aligned} \label{matrix}
\end{equation}
Both eigenvalues are positive:
\begin{equation}
 w+x=4C^2\lambda(1-q)+\frac{32C^2\lambda^2q}{1+q},\qquad
 w-x=4C^2\lambda(1+q).
\end{equation}
The symmetric and antisymmetric force combinations consequently
have the same normalized two-point kernel but different strengths.

\newpage
\section{The projected deficit and its cross with the quartic field}
Put $v=2\lambda(1-q)+16\lambda^2q/(1+q)$, $d=2\lambda(1+q)$,
and $\rho=v/(v+d)$. The preceding equal-time Gram and its
nondegenerate finite-$N$ limit give, for
$\Pi_N=\operatorname{proj}_{\{I,B,O_-,J_hF\}}$,
\begin{equation}
 T_N=(1-\Pi_N)\mathcal S_N
 =\frac{Y_-+Y_+}{2}-\rho Y_-+o_{L^2}(1). \label{T}
\end{equation}
The coefficient of the additional $J_hF$ direction tends to zero;
its norm is bounded and its limiting Gram is nondegenerate.
No unproved energy-resolved $F$ cross is needed for this replacement.
Equation~\eqref{matrix} now gives
\begin{equation}
 d\nu_T=V_{\rm new}\,d\mu_{q^2},\qquad
 V_{\rm new}=\frac{C^2vd}{v+d}
 =\frac{C^2\lambda(1+q)(1-q^2+8\lambda q)}{1+q+4\lambda q}.
\end{equation}

The remaining issue is the cross with $Q_4$. To make the finite
certificate reproducible, define
\begin{equation}
 d_3=q^2-q+1,\quad d_5=1+q+q^2+q^3+q^4,\quad D_4=(1+q^2)d_3d_5.
\end{equation}
In the four-slot mark order $01,02,03,12,13,23$, the numerator
weights for $Q_4$ are
\begin{align}
 \boldsymbol w_4=\big(&q^3(1+q^2),-q(1+q^2)d_5,D_4,\notag\\
 &q^2(1+q+q^2)(1+q^2),-q(1+q^2)d_5,q^3(1+q^2)\big).
\end{align}
The unmarked numerator is
$2\lambda(q-1)^3(1+q)(1+q+q^2)d_3$.
For $Y_+$, the marked and unmarked numerators are
$1+q$ and $-2\lambda(1-q)$, with denominator $1+q$.

Pair a four-slot block with a two-slot block. If they share
$k=0,1,2$ bridges, the remaining $6-2k$ endpoints have an $A/B$
color word and one split between the two exterior arcs. There are
\begin{equation}
 \sum_{k=0}^2\binom{6-2k}{4-k}(7-2k)=128
\end{equation}
templates, each with $4!2!/k!$ pairings. Substituting the
opposite-orientation rule~\eqref{minor}, all 128 coefficients of
$\lambda$ and $\lambda^2$ vanish as integer polynomials in $q$.
The certificate retains all 6072 words' template aggregates.
Arbitrary spectator chords contribute a common factor on each
endpoint orbit: continuous block permutations preserve crossings
with all exterior chords except the explicitly retained endpoints.
Thus this finite classification proves cancellation for arbitrary
fixed two-arc energy degrees, not just the nine low moments checked
separately. Together with the old normal-orientation certificate,
\begin{equation}
 \nu_{Q_4,Y_-}=\nu_{Q_4,Y_+}=0,\qquad \nu_{Q_4,T}=0. \label{zero}
\end{equation}
Equation~\eqref{matrix} also gives $\nu_{T,Y_-}=0$ at every energy,
not only after energy integration. Thus $\nu_{X,Y_-}=0$, including
crosses with all fixed left/right energy descendants of $Y_-$.

\newpage
\section{From moments to the full memory spectrum}
For Hermitian fields $X_i$, the finite-system matrix measure is
\begin{equation}
 \nu_{ij,N}(dE,dE')=
 \EE\tr_m\big[X_i\Pi_h(dE)X_j\Pi_h(dE')\big].
\end{equation}
It is a positive matrix-valued measure, as seen in an energy
eigenbasis. The fixed-word overlap bounds give all diagonal energy
moments, tightness, and uniform integrability of any fixed moment.
Spectral Cauchy--Schwarz bounds every off-diagonal total variation
by the diagonal measures. The candidate measures have compact
energy support, so their full moment sequence is determinate.
Polarization therefore gives joint weak convergence for
$(Q_4,Y_-,Y_+)$. In particular~\eqref{zero} is a statement about
the entire cross measure, not merely the equal-time inner product.

An $L^2$ replacement $X-X'$ changes the self measure by total
variation at most $\|X-X'\|_2(\|X\|_2+\|X'\|_2)$, by spectral
Cauchy--Schwarz on partitions. Hence~\eqref{strong} and~\eqref{T}
also apply to bounded spectral functions without interchanging an
unbounded exponential series and the limit.

Ref.~\cite{desc} established the actual raw-coefficient reduction
\begin{equation}
 (1-\widehat\Pi_N)mZ_2
 =Q_{4,N}+8\lambda_N C_Nq_N^2T_N+o_{L^2}(1).
\end{equation}
The old fourth-chaos theorem and~\eqref{zero} prove~\eqref{result},
with
\begin{align}
 V_2&=64\J^8q^4
 \frac{(1+q)(1-q^2+8\lambda q)}{\lambda(1+q+4\lambda q)},\\
 V_4&=C^4\frac{4\lambda(1-q)^2(1+q)^2
 [P(q)+4\lambda q^3 U(q)]}{(1+q^2)d_3d_5},\\
 P(q)&=1+q+2q^2+q^3+q^4,\qquad
 U(q)=2+q+2q^2+q^3+2q^4.
\end{align}
For completeness, the existing probe kernel~\cite{TT,BNR} is
\begin{align}
 d\mu_r&=\rho_q(\theta)\rho_q(\phi)K_r(\theta,\phi)
                      \,d\theta\,d\phi,\\
 \rho_q(\theta)&=\frac{(q;q)_\infty
                  |(e^{2i\theta};q)_\infty|^2}{2\pi},\\
 K_r(\theta,\phi)&=\frac{(r^2;q)_\infty}
 {\prod_{s,t=\pm1}(re^{i(s\theta+t\phi)};q)_\infty},\\
 E(\theta)&=\frac{2\sqrt C\cos\theta}{\sqrt{1-q}},\qquad
 G_r(t)=\int e^{it(E-E')}\,d\mu_r.
\end{align}
Here $\theta,\phi\in[0,\pi]$ and $(z;q)_\infty=\prod_{j\geq0}(1-zq^j)$.
For fixed $0<q<1$ the measures are positive and absolutely continuous;
there is no limiting static atom. Their use here follows the
same-coupling composite-field reduction, rather than assuming an
independent random probe at the outset.

\newpage
\section{Small deformation, evidence, and continuation}
After the double-scaled limit, take $\lambda\downarrow0$ at fixed
string time. The known probe limit gives
\begin{align}
 G_{q^2}(t)&\longrightarrow\operatorname{sech}^4(2\J t),&
 G_{q^4}(t)&\longrightarrow\operatorname{sech}^8(2\J t),\\
 V_2&\longrightarrow1024\J^8,&
 V_4&\sim\frac{768\J^8}{5\lambda},\\
 \frac{V_2}{V_4+V_2}&=\frac{20}{3}\lambda+o(\lambda),&
 \frac{C_X(t)}{V_4+V_2}&\longrightarrow\operatorname{sech}^8(2\J t).
\end{align}
Thus the lower-chaos contribution survives with finite absolute
strength while its fraction vanishes in this ordered limit.
The coefficient $20/3$ describes the weight fraction only; it is
not the complete first correction to the time kernel, which also
receives the deformation of $G_{q^4}$.

\paragraph{Executed checks.}
The local normal/reverse determinant rule was checked in 1688
four-, six-, and eight-slot cases. Direct CAR matrices supplied
640 words and 2560 comparisons, including 115 nonzero double-mark
values and three negative values. Five small half-filled systems,
each with three deterministic nonzero Hermitian samples, supplied
90 exact bridge identities. Cleared complex-integer entries stayed
below $2^{48}$; the Casimir implementation uses $2\mathcal C_x$.

The new opposite-orientation endpoint certificate has 128 templates
and no nonzero polynomial. A separate direct pairing enumeration
checks nine $(a,b)$ mixed moments through $a+b=4$, including the
single-mark and double-mark sums. Symbolic checks reproduce the
previous deficit Gram and the Schur residual $V_{\rm new}$, both
positive eigenvalues, and the two small-$\lambda$ coefficients.

For $\J=1$, $\lambda=0.25,0.1,0.025$ and
$t=0,0.25,0.5,1$, a symmetric $q$-Jacobi evolution evaluates both
kernels and their mixture. A Taylor bound controls normalized
truncation below $10^{-13}$; twelve extra Jacobi levels check
refinement. At the first two couplings, independent 128- and
192-node Mehler quadrature agrees within $2\times10^{-10}$.
The numerical file retains all weights, values, cutoffs, and errors.
Floating rounding and quadrature are refinement checks, not
interval certificates. Source hashes and dependency versions are
stored alongside each result.

\paragraph{Next physical step.}
The specified raw $mZ_2$ projection is now closed at two points.
The next tasks are its mixed spectra with earlier memory fields,
a uniform central-charge-window extension, and a finite-memory
two-parameter remainder. These feed the Susskind singlet/meson
program. A chord label $2$ or $4$ alone does not fix a bulk meson
mass or a spatial reconstruction map.

\begin{thebibliography}{9}\small
\bibitem{TT} J. Tall and S. Tomsovic,
\emph{Double-scaled bosonic and fermionic embedded ensembles,
complex SYK, and the dual Hilbert space},
\href{https://arxiv.org/html/2604.14522v2}{arXiv:2604.14522v2},
Secs.~3.1 and 4.1. Spectral $q$-Wick products and existing strength kernels.
\bibitem{BNR} M. Berkooz, V. Narovlansky and H. Raj,
\href{https://arxiv.org/abs/2006.13983}{arXiv:2006.13983}.
Complex-SYK chord and probe conventions, as recorded in the project source notes.
\bibitem{desc} A. Day and Chengheng, project notes:
\emph{Restoring the raw memory descendants} (2026-11-06);
\emph{Second memory coefficient and fourth singlet};
\emph{From derivative singlet to its full two-point spectrum}.
Local derivations and machine-readable certificates are preserved in the project.
\end{thebibliography}
\end{document}
